% Figure: geometric comparison of the Wald, score, and
% likelihood-ratio test statistics on a single log-likelihood curve.
% The curve is a generic concave log-likelihood ell(theta) plotted
% against theta, with the MLE theta-hat at the peak and a null value
% theta_0 to its left. The three statistics read off three different
% geometric features: the LR test from the vertical drop, the Wald
% test from the horizontal distance scaled by curvature, the score
% test from the slope at theta_0. Coordinates are a precomputed
% quadratic-with-cubic log-likelihood for illustration.
\begin{figure}[htbp]
\centering
\begin{tikzpicture}
\begin{axis}[
    width=12cm, height=7cm,
    xmin=-0.2, xmax=3.6,
    ymin=-4.6, ymax=0.6,
    axis lines=left,
    xlabel={Parameter $\theta$},
    ylabel={Log-likelihood $\ell(\theta)$},
    every axis plot/.append style={very thick, mark=none},
    legend style={font=\small, draw=none, fill=none, at={(0.98,0.05)}, anchor=south east},
    xtick={0.6, 2.4},
    xticklabels={$\theta_0$, $\hat{\theta}$},
    ytick=\empty,
]
% Log-likelihood curve: a smooth concave function peaking at 2.4.
\addplot[engblack, domain=-0.1:3.5, samples=80]{-0.55*(x-2.4)^2 + 0.03*(x-2.4)^3};
\addlegendentry{$\ell(\theta)$}
% Vertical drop = likelihood-ratio statistic (-2[ell(theta0)-ell(hat)]).
\addplot[engred, dashed] coordinates {(0.6, 0.0) (0.6, -1.79)};
% ell at theta0 = -0.55*(0.6-2.4)^2 + 0.03*(0.6-2.4)^3 = -1.783 - 0.175 = -1.958; round to label
% Tangent line at theta0 (slope = score) for the score test.
\addplot[engblue, domain=0.0:1.4, samples=2]{-1.96 + 1.815*(x-0.6)};
\addlegendentry{tangent at $\theta_0$ (score)}
% Vertical line at MLE.
\addplot[enggray!50, dotted] coordinates {(2.4, -4.6) (2.4, 0.0)};
% Vertical line at theta0.
\addplot[enggray!50, dotted] coordinates {(0.6, -4.6) (0.6, -1.96)};
% Points.
\addplot[only marks, engblack, mark=*, mark size=2pt] coordinates {(2.4, 0.0) (0.6, -1.96)};
% Horizontal brace marker for Wald (distance theta-hat - theta0).
\addplot[englime!70!black, very thick] coordinates {(0.6, -4.3) (2.4, -4.3)};
\node[anchor=north, font=\footnotesize, englime!60!black] at (axis cs:1.5, -4.3) {Wald: $\hat{\theta} - \theta_0$, scaled by curvature};
\node[anchor=west, font=\footnotesize, engred] at (axis cs:0.66, -0.95) {LR drop};
\end{axis}
\end{tikzpicture}
\caption[Geometry of the Wald, score, and likelihood-ratio
tests]{Three classical tests of $H_0: \theta = \theta_0$ read three
different features off the same log-likelihood curve. The
likelihood-ratio statistic is twice the vertical drop from the peak
$\ell(\hat{\theta})$ to $\ell(\theta_0)$. The Wald statistic uses the
horizontal distance $\hat{\theta} - \theta_0$, scaled by the
curvature at the maximum. The score statistic uses the slope of the
tangent at $\theta_0$, scaled by the information there. The three
agree to leading order under the null and disagree away from it; the
likelihood-ratio test, invariant under reparametrisation, is the one
to prefer when they diverge.}
\label{fig:vol02:ch14:trinity-tests}
\end{figure}
